\documentclass[10pt,a4paper]{article}
\usepackage{../../styles/neet_style}

\pagestyle{fancy}
\fancyhf{}
\fancyhead[L]{\small \textbf{NEET UG 2027 Master Notes}}
\fancyhead[R]{\small \textbf{Satya Prakash}}
\fancyfoot[L]{\small \textbf{Zoology Chapter Series}}
\fancyfoot[C]{\thepage}
\fancyfoot[R]{\small \textbf{Human Health \& Evolution}}

\begin{document}

\begin{center}
  {\LARGE \textbf{HUMAN HEALTH, IMMUNOLOGY \& EVOLUTIONARY BIOLOGY}}\\[4pt]
  {\large \textbf{Comprehensive High-Yield Master Notes: NMC Rationalized Core Modules}}\\[6pt]
  \textbf{Prepared, Compiled \& Authored by Satya Prakash}
\end{center}

\vspace{-4pt}
\hrule height 1pt
\vspace{8pt}

\section{Principles of Immunology \& Host Defense Mechanisms}

\begin{conceptbox}[Innate Immunity: Non-Specific First \& Second Lines of Defense]
Present from birth, non-specific, lacks immunological memory. Comprises four distinct anatomical and physiological barrier tiers:
\begin{enumerate}[leftmargin=*]
  \item \textbf{Physical Barriers:} Skin (keratinized stratum corneum prevents microbial ingress); Mucous coating of epithelial lining across respiratory, gastrointestinal, and urogenital tracts traps entrapped microbes.
  \item \textbf{Physiological Barriers:} Gastric hydrochloric acid ($\mathrm{pH} \approx 1.5-2.0$ destroys ingested bacteria); Saliva and Tears contain \textbf{lysozyme} (hydrolyzes $\beta(1\to 4)$ glycosidic bonds in bacterial peptidoglycan cell walls).
  \item \textbf{Cellular Barriers:} Phagocytes that engulf and digest foreign pathogens:
  \begin{itemize}
    \item \textbf{PMNL} (Polymorphonuclear Leukocytes / Neutrophils): Most abundant circulating WBCs ($60-65\%$).
    \item \textbf{Monocytes} ($6-8\%$) differentiating into tissue \textbf{Macrophages} (e.g., Kupffer cells, microglia, alveolar macrophages).
    \item \textbf{Natural Killer (NK) Cells:} Specialized cytotoxic lymphocytes that destroy virus-infected and malignant tumor cells non-specifically via perforins and granzymes.
  \end{itemize}
  \item \textbf{Cytokine Barriers:} Virus-infected cells secrete non-viral glycoproteins called \textbf{Interferons ($\alpha, \beta, \gamma$)}, which bind to uninfected neighboring cells, inducing translation of antiviral proteins (AVPs) that inhibit viral translation and replication.
\end{enumerate}
\end{conceptbox}

\begin{conceptbox}[Acquired (Adaptive) Immunity: Humoral vs Cell-Mediated Immunity (CMI)]
Characterized by antigenic specificity, diversity, discrimination between self and non-self, and immunological memory.
\begin{center}
\scriptsize
\begin{tabular}{|p{2.4cm}|p{6.2cm}|p{6.2cm}|}
\hline
\textbf{Property} & \textbf{Humoral / Antibody-Mediated Immunity (AMI)} & \textbf{Cell-Mediated Immunity (CMI)} \\
\hline
\textbf{Effector Cells} & \textbf{B-lymphocytes} differentiating into Plasma cells (antibody factories) and Memory B cells. & \textbf{T-lymphocytes} maturing into Helper T ($\mathrm{CD4}^+ \mathrm{T_H}$), Cytotoxic T ($\mathrm{CD8}^+ \mathrm{T_C}$), and Regulatory T cells. \\
\hline
\textbf{Mode of Action} & Secretes soluble \textbf{Antibodies} ($\mathrm{H_2L_2}$) into blood plasma, lymph, and extracellular fluids. & Direct cell-to-cell contact; $\mathrm{T_C}$ cells release perforins creating pores in target membranes; activates macrophages. \\
\hline
\textbf{Target Pathogens} & Extracellular bacteria, viruses in circulation, circulating venoms, toxins. & \textbf{Intracellular pathogens} (\textit{Mycobacterium tuberculosis}, viruses inside cells, fungi, parasites) and cancerous cells. \\
\hline
\textbf{Organ Graft Rejection} & Plays a minor secondary role (hyperacute humoral rejection). & \textbf{Primary Mediator of Graft Rejection!} Allograft rejection is mediated by $\mathrm{T_C}$ cells recognizing non-self MHC/HLA antigens. \\
\hline
\textbf{Immunosuppression} & Intravenous immunoglobulins (IVIg), plasmapheresis. & \textbf{Cyclosporin A} (fungal metabolite from \textit{Trichoderma polysporum}) inhibits calcineurin, suppressing T-cell activation. \\
\hline
\end{tabular}
\end{center}
\end{conceptbox}

\begin{formulabox}[Antibody Architecture ($\mathrm{H_2L_2}$) \& Immunoglobulin Classes]
Each antibody is a Y-shaped heterodimer composed of \textbf{4 polypeptide chains}: two identical Heavy ($\mathrm{H}$) chains ($\approx 50\,\mathrm{kDa}$ each) and two identical Light ($\mathrm{L}$) chains ($\approx 25\,\mathrm{kDa}$ each), cross-linked by \textbf{interchain disulfide bonds} ($\mathrm{-S-S-}$).
\begin{equation}
  \mathbf{\text{Antibody Formula: } \mathrm{H_2 L_2}} \quad (\text{Two Heavy Chains } \mathrm{H} + \text{Two Light Chains } \mathrm{L})
\end{equation}
\centerline{\small Fab (Antigen-Binding Fragment) contains Variable ($V_H, V_L$) domain; Fc mediates effector binding.}
\vspace{2pt}
\begin{center}
\scriptsize
\begin{tabular}{|p{1.2cm}|p{2.8cm}|p{10.6cm}|}
\hline
\textbf{Class} & \textbf{Valency / Form} & \textbf{Diagnostic Role \& High-Yield NCERT Facts} \\
\hline
$\mathbf{\mathrm{IgG}}$ & Monomer (Bivalent) & Most abundant ($75-80\%$ of serum Igs); \textbf{only antibody class that crosses human placenta} providing natural passive immunity to fetus. \\
\hline
$\mathbf{\mathrm{IgA}}$ & Dimer (Tetravalent) & Secretory Ig linked by J-chain and secretory piece. Abundant in \textbf{Colostrum} (first yellowish milk), saliva, tears, respiratory mucous. \\
\hline
$\mathbf{\mathrm{IgM}}$ & Pentamer (Decavalent) & Largest molecular weight, linked by J-chain; first antibody synthesized during primary immune response; potent complement activator. \\
\hline
$\mathbf{\mathrm{IgE}}$ & Monomer (Bivalent) & Mediates \textbf{Type I Hypersensitivity (Allergies)}; binds Fc receptors on mast cells and basophils, triggering release of histamine and serotonin. \\
\hline
$\mathbf{\mathrm{IgD}}$ & Monomer (Bivalent) & Co-expressed with $\mathrm{IgM}$ on the surface of naive mature B-cells; functions as antigen receptor initiating B-cell differentiation. \\
\hline
\end{tabular}
\end{center}
\end{formulabox}

\begin{conceptbox}[Lymphoid Organs: Primary vs Secondary Tier Architecture]
\begin{itemize}[leftmargin=*]
  \item \textbf{Primary Lymphoid Organs (Generative/Central):} Sites where immature lymphocytes differentiate and proliferate into antigen-sensitive lymphocytes.
  \begin{itemize}
    \item \textbf{Bone Marrow:} Primary site of hematopoiesis; produces all blood cells including B- and T-progenitors. B-cells mature completely here.
    \item \textbf{Thymus:} Lobed organ situated near heart beneath breastbone; large at birth, undergoes progressive atrophy with age. Provides microenvironment for \textbf{T-lymphocyte maturation and thymic selection}.
  \end{itemize}
  \item \textbf{Secondary Lymphoid Organs (Peripheral):} Sites where lymphocytes interact with antigens and proliferate into effector/memory cells.
  \begin{itemize}
    \item \textbf{Spleen:} Large bean-shaped organ containing lymphocytes and phagocytes. Acts as a physiological blood filter trapping blood-borne microbes; serves as a massive \textbf{reservoir of erythrocytes (RBCs)}.
    \item \textbf{Lymph Nodes:} Small solid structures along lymphatic vessels; filter lymph fluid by trapping micro-organisms, presenting them to resident dendritic cells and lymphocytes.
    \item \textbf{MALT (Mucosa-Associated Lymphoid Tissue):} Located within mucosal lining of major tracts (digestive, respiratory, urogenital). Constitutes \textbf{$\approx 50\%$ of all lymphoid tissue in the human body!} Includes Peyer's patches (ileum) and palatine tonsils.
  \end{itemize}
\end{itemize}
\end{conceptbox}

\section{Major Human Pathogens, Transmission \& Diagnostic Tests}

\begin{conceptbox}[Bacterial, Viral, Protozoan \& Helminthic Infectious Diseases]
\begin{center}
\scriptsize
\begin{tabular}{|p{1.8cm}|p{2.4cm}|p{2.4cm}|p{7.6cm}|}
\hline
\textbf{Disease} & \textbf{Causative Organism} & \textbf{Transmission Vector} & \textbf{Diagnostic Pathology \& NCERT Key Hallmarks} \\
\hline
\textbf{Typhoid} & \textit{Salmonella typhi} (Gram-negative bacillus) & Contaminated food and water; classic carrier: \textbf{Mary Mallon} (Typhoid Mary). & Sustained high fever ($39-40^\circ\mathrm{C}$), weakness, stomach pain, constipation, headache, anorexia. \textbf{Intestinal perforation} in severe cases. Confirmatory test: \textbf{Widal Test}. \\
\hline
\textbf{Pneumonia} & \textit{Streptococcus pneumoniae}, \textit{Haemophilus influenzae} & Inhaling droplets/aerosols; sharing glasses/utensils. & \textbf{Infects pulmonary alveoli}; alveoli become fluid-filled leading to severe respiratory distress. Lips and fingernails turn \textbf{gray to bluish} in severe cases. \\
\hline
\textbf{Common Cold} & \textbf{Rhino viruses} (ssRNA) & Inhaling droplets; contaminated pens, books, cups. & \textbf{Infects nose and respiratory passage, but NOT lungs!} Nasal congestion, discharge, sore throat, hoarseness, cough; lasts $3-7$ days. \\
\hline
\textbf{Amoebiasis} & \textit{Entamoeba histolytica} (Protozoan parasite) & \textbf{Houseflies} act as mechanical carriers from feces to food. & Infects large intestine; constipation, abdominal pain, cramps, stools with \textbf{excess mucous and blood clots}. \\
\hline
\textbf{Ascariasis} & \textit{Ascaris lumbricoides} (Intestinal roundworm) & Ingestion of water, vegetables, fruits contaminated with fecal eggs. & Internal bleeding, muscular pain, fever, anemia, and mechanical \textbf{blockage of the intestinal passage}. \\
\hline
\textbf{Filariasis (Elephantiasis)} & \textit{Wuchereria bancrofti}, \textit{W. malayi} (Filarial worms) & Bite of infected \textbf{female \textit{Culex}} mosquito. & Chronic inflammation of organs, especially \textbf{lymphatic vessels of lower limbs and genital organs}, resulting in gross gross deformities. \\
\hline
\textbf{Ringworm} & \textit{Microsporum}, \textit{Trichophyton}, \textit{Epidermophyton} & Direct contact, sharing towels, clothes, comb; heat and moisture favor growth. & Dry, scaly, pruritic (itchy) circular lesions on skin, nails, scalp, and groin folds. \\
\hline
\end{tabular}
\end{center}
\end{conceptbox}

\begin{conceptbox}[Life Cycle of \textit{Plasmodium} (Malarial Parasite: Digenetic Cycle)]
\textit{Plasmodium} requires two hosts: \textbf{Primary host} = Female \textit{Anopheles} mosquito (sexual cycle); \textbf{Secondary host} = Human (asexual cycle).
\begin{enumerate}[leftmargin=*]
  \item \textbf{Inoculation:} Female \textit{Anopheles} mosquito injects infective stage \textbf{Sporozoites} with saliva into human blood during feeding.
  \item \textbf{Hepatic Schizogony:} Sporozoites invade liver parenchymal cells, multiply asexually, and burst hepatocytes, releasing merozoites.
  \item \textbf{Erythrocytic Schizogony:} Merozoites penetrate RBCs, form ring-stage trophozoites $\to$ schizonts, and burst RBCs synchronously every $48-72$ hours.
  \begin{itemize}
    \item Rupture of RBCs releases pyrogenic toxin \textbf{Haemozoin} (formed from digested hemoglobin iron-porphyrin), producing \textbf{chills and high recurring fever}.
  \end{itemize}
  \item \textbf{Gametocytogenesis:} Some merozoites in RBCs develop into sexual stages: male microgametocytes and female macrogametocytes.
  \item \textbf{Mosquito Phase (Sexual Cycle):} Female \textit{Anopheles} ingests gametocytes with blood meal. Fertilization and zygote formation occur in the \textbf{mosquito gut}. Zygote becomes motile \textbf{ookinete}, penetrates gut wall to form \textbf{oocyst}. Oocyst divides by sporogony, generating thousands of sporozoites which migrate to and lodge in the \textbf{mosquito salivary glands}.
\end{enumerate}
\end{conceptbox}

\begin{conceptbox}[HIV/AIDS Pathogenesis \& Viral Replication Mechanism]
Caused by \textbf{Human Immunodeficiency Virus (HIV)}, an enveloped Lentivirus (Retroviridae) harboring \textbf{two identical single-stranded RNA (+ssRNA)} molecules and \textbf{Reverse Transcriptase}.
\begin{enumerate}[leftmargin=*]
  \item \textbf{Host Cell Attachment:} Viral envelope glycoprotein $\mathbf{\mathrm{gp120}}$ binds to \textbf{$\mathrm{CD4}$ receptors} and chemokine coreceptors ($\mathrm{CCR5/CXCR4}$) on macrophages and Helper T ($\mathrm{T_H}$) lymphocytes.
  \item \textbf{Macrophage Phase (``HIV Factory''):} Viral RNA is converted into viral DNA by Reverse Transcriptase ($\text{RNA} \to \text{cDNA} \to \text{dsDNA}$). Viral DNA integrates into host genome via \textbf{Integrase}. Macrophages produce viral particles without immediate cell lysis.
  \item \textbf{$\mathrm{T_H}$ Cell Depletion:} Virus enters $\mathrm{CD4}^+ \mathrm{T_H}$ cells, replicates, and lyses host cells. Helper T-cell count plummets from normal ($900-1200/\mu\mathrm{L}$) to below $\mathbf{200/\mu\mathrm{L}}$.
  \item \textbf{Opportunistic Infections:} Severe cell-mediated immunodeficiency permits life-threatening infections by \textit{Mycobacterium tuberculosis}, \textit{Toxoplasma gondii}, \textit{Pneumocystis jirovecii}, and cytomegalovirus (CMV).
  \item \textbf{Diagnosis \& Management:} Screening test: \textbf{ELISA} (Enzyme-Linked Immunosorbent Assay); Confirmatory test: \textbf{Western Blot} (detects viral antigens $\mathrm{p24, gp120, gp41}$). Treatment: Anti-Retroviral Therapy (\textbf{ART}) cocktail of nucleoside reverse transcriptase inhibitors (NRTIs, e.g., Zidovudine/AZT) and protease inhibitors.
\end{enumerate}
\end{conceptbox}

\begin{conceptbox}[Cancer Biology: Oncogenesis, Characteristics \& Clinical Therapies]
\begin{itemize}[leftmargin=*]
  \item \textbf{Transformation Hallmarks:} Normal cells exhibit \textbf{contact inhibition} (contact with neighboring cells inhibits uncontrolled proliferation). Cancer cells lose contact inhibition, proliferating relentlessly into \textbf{tumors}:
  \begin{itemize}
    \item \textbf{Benign Tumors:} Encapsulated, remain localized to primary site, do not invade adjacent tissues.
    \item \textbf{Malignant Tumors:} Unencapsulated masses of neoplastic cells exhibiting \textbf{Metastasis} (cells slough off into circulation and colonize distant organs; \textbf{the most feared property of malignant neoplasms}).
  \end{itemize}
  \item \textbf{Etiological Carcinogens:}
  \begin{itemize}
    \item \textbf{Physical:} Ionizing radiation (X-rays, $\gamma$-rays break DNA strands); Non-ionizing radiation (\textbf{UV-B} induces cyclobutane pyrimidine/thymine dimers).
    \item \textbf{Chemical:} Tobacco smoke contains \textbf{$N$-nitrosamines} and benzo[$a$]pyrene (major cause of bronchogenic lung carcinoma).
    \item \textbf{Biological:} Oncogenic viruses carrying viral oncogenes ($v\text{-}onc$); cellular proto-oncogenes ($c\text{-}onc$) mutated into active oncogenes.
  \end{itemize}
  \item \textbf{Therapeutic Modalities:} Surgery (tumor excision); Radiotherapy (targeted lethal irradiation of malignant tissue); Chemotherapy (antimitotic drugs like Vincristine, Vinblastine from \textit{Catharanthus roseus}); \textbf{Immunotherapy:} Biological response modifiers like \textbf{$\alpha$-Interferon} activate the patient's immune system to destroy tumors.
\end{itemize}
\end{conceptbox}

\begin{conceptbox}[Psychoactive Drugs, Hallucinogens \& Addiction Pharmacology]
\begin{center}
\scriptsize
\begin{tabular}{|p{1.8cm}|p{2.5cm}|p{2.5cm}|p{7.5cm}|}
\hline
\textbf{Class} & \textbf{Representative Drugs} & \textbf{Receptors / Site} & \textbf{Pharmacological Action \& Botanical Extraction} \\
\hline
\textbf{Opioids} & Morphine, Heroin (\textbf{Smack} / Diacetylmorphine), Codeine & Specific opioid receptors in \textbf{CNS and GIT}. & Extracted from latex of unripe capsules of \textbf{\textit{Papaver somniferum}} (Opium poppy). Heroin is synthesized by \textbf{acetylation of morphine}; white, odorless, bitter crystalline compound. Acts as a \textbf{depressant and slows down body functions}. \\
\hline
\textbf{Cannabinoids} & Marijuana, Hashish, Charas, Ganja & Cannabinoid receptors primarily located in the \textbf{brain}. & Obtained from inflorescences, leaves, and resin of \textbf{\textit{Cannabis sativa}} (Hemp plant). Inhaled or orally ingested; induces relaxation, alters sensory perception, and severely impacts \textbf{cardiovascular system}. \\
\hline
\textbf{Coca Alkaloids} & \textbf{Cocaine} (Coke / Crack) & Interferes with reuptake transport of \textbf{Dopamine}. & Extracted from leaves of South American shrub \textbf{\textit{Erythroxylum coca}}. Potent CNS stimulant producing euphoria and increased energy; high doses induce \textbf{hallucinations}. Other hallucinogens: \textit{Atropa belladonna}, \textit{Datura}. \\
\hline
\textbf{Tobacco / Nicotine} & Nicotine (Alkaloid) & Stimulates nicotinic acetylcholine receptors. & Stimulates \textbf{adrenal medulla} to release adrenaline and noradrenaline into circulation, raising blood pressure and heart rate. Induces lung cancer, emphysema, bronchitis; elevates carbon monoxide ($\mathrm{CO}$) in blood, decreasing oxyhemoglobin concentration. \\
\hline
\end{tabular}
\end{center}
\end{conceptbox}

\section{Evolutionary Biology: Origins, Evidence \& Mechanisms}

\begin{conceptbox}[Origin of Life: Abiogenesis \& Miller-Urey Experimental Proof]
\begin{itemize}[leftmargin=*]
  \item \textbf{Primordial Earth Atmosphere:} Formed $\approx 4.5$ billion years ago. Primitive atmosphere was strictly \textbf{reducing}: no free $\mathrm{O_2}$, rich in $\mathrm{CH_4, NH_3, H_2O\text{ (vapor)}, H_2}$.
  \item \textbf{Chemical Evolution (Oparin-Haldane Hypothesis):} First forms of life originated from pre-existing non-living organic molecules (RNA, proteins, sugars), preceding biological evolution.
  \item \textbf{Miller-Urey Spark-Discharge Experiment (1953):} Stanley Miller simulated prebiotic atmospheric conditions:
  \begin{itemize}
    \item Closed spark-discharge flask containing gases: $\mathbf{\mathrm{CH_4 : NH_3 : H_2}}$ in ratio $\mathbf{2 : 1 : 2}$ along with continuously boiling water vapor ($\mathrm{H_2O}$).
    \item Tungsten electrodes produced continuous electrical discharge ($\approx 75,000\,\mathrm{V}$) maintaining chamber temperature at $\mathbf{800^\circ\mathrm{C}}$.
    \item \textbf{Experimental Yield:} Condensed effluent contained amino acids: \textbf{Glycine, $\alpha$-Alanine, and Aspartic acid}, conclusively validating chemical evolution.
  \end{itemize}
\end{itemize}
\end{conceptbox}

\begin{conceptbox}[Comparative Morphological Evidences: Homology vs Analogy]
\begin{center}
\scriptsize
\begin{tabular}{|p{2.5cm}|p{6.0cm}|p{6.0cm}|}
\hline
\textbf{Parameter} & \textbf{Homologous Organs (Divergent Evolution)} & \textbf{Analogous Organs (Convergent Evolution)} \\
\hline
\textbf{Fundamental Basis} & \textbf{Same anatomical origin and basic structure}, modified to perform different functions in response to different adaptations. & \textbf{Different anatomical origin and basic structure}, modified to perform similar functions due to similar selective pressures. \\
\hline
\textbf{Evolutionary Pattern} & \textbf{Divergent Evolution} (Radiates from a common ancestor into diverse ecological niches). & \textbf{Convergent Evolution} (Unrelated lineages evolve functionally convergent phenotypes). \\
\hline
\textbf{Vertebrate Forelimb Examples} & Forelimb skeletal architecture of \textbf{Whale (flipper), Bat (patagium), Cheetah (running), and Human (grasping)} shares identical arrangement: humerus, radius, ulna, carpals, metacarpals, phalanges. & \textbf{Wings of Butterfly} (chitinous membranous folds) vs \textbf{Wings of Birds} (feather-covered forelimb bones). \\
\hline
\textbf{Sensory / Functional Organs} & Vertebrate heart and brain anatomy across fish, amphibians, reptiles, birds, and mammals. & \textbf{Eye of Octopus} (invertebrate retina derived from skin) vs \textbf{Eye of Mammals} (derived from embryonic diencephalon). \textbf{Flippers of Penguins (bird)} vs \textbf{Dolphins (mammal)}. \\
\hline
\textbf{Botanical Examples} & \textbf{Thorns of \textit{Bougainvillea}} (axillary bud protection) vs \textbf{Tendrils of \textit{Cucurbita}} (axillary bud climbing). & \textbf{Sweet Potato} (root tuber modification) vs \textbf{Potato} (underground stem tuber modification). \\
\hline
\end{tabular}
\end{center}
\end{conceptbox}

\begin{conceptbox}[Adaptive Radiation \& Convergent Marsupial-Placental Counterparts]
\textbf{Adaptive Radiation:} Evolution of divergent species starting from a common ancestral point and radiating into distinct geographical habitats and ecological niches.
\begin{itemize}[leftmargin=*]
  \item \textbf{Classic Examples:} Darwin's Finches on Galapagos archipelago (ancestral seed-eating beak diversified into vegetarian, insectivorous, warbler, and cactus-puncturing beaks); Australian Marsupials radiating from a single ancestral stock within the isolated island continent.
  \item \textbf{Marsupial vs Placental Mammal Convergence:} When adaptive radiation occurs in isolated continents for different lineages occupying similar ecological niches, convergent evolution generates striking morphological analogs:
  \begin{center}
  \scriptsize
  \begin{tabular}{|l|l|l|}
  \hline
  \textbf{Ecological Niche} & \textbf{Australian Marsupial Counterpart} & \textbf{Placental Mammal Counterpart} \\
  \hline
  Fossorial Burrower & Marsupial mole (\textit{Notoryctes}) & Placental Mole \\
  Myrmecophage (Ant-eater) & Numbat / Banded anteater (\textit{Myrmecobius}) & Anteater \\
  Small Glider & Flying phalanger / Sugar glider (\textit{Petaurus}) & Flying squirrel \\
  Feline Carnivore & Tasmanian tiger cat (\textit{Dasyurus}) & Bobcat \\
  Apex Carnivore & Tasmanian wolf (\textit{Thylacinus}) & Wolf (\textit{Canis lupus}) \\
  \hline
  \end{tabular}
  \end{center}
\end{itemize}
\end{conceptbox}

\begin{formulabox}[Hardy-Weinberg Equilibrium: Mathematical Formulations \& Selection Dynamics]
In a large, randomly mating, panmictic diploid population with no evolutionary disruptive forces, allele frequencies and genotype frequencies remain constant across generations:
\begin{equation}
  \mathbf{p + q = 1} \qquad \implies \qquad \mathbf{p^2 + 2pq + q^2 = 1}
\end{equation}
Where $p = \text{frequency of dominant allele } A$, $q = \text{frequency of recessive allele } a$, $p^2 = \text{frequency of homozygous dominant } AA$, $2pq = \text{heterozygous } Aa$, and $q^2 = \text{homozygous recessive } aa$.
\begin{itemize}[leftmargin=*]
  \item \textbf{Five Evolutionary Forces Disrupting Equilibrium:} (1) Gene migration / Gene flow; (2) \textbf{Genetic Drift} (random allele frequency fluctuation in small populations, including \textbf{Founder Effect} and \textbf{Bottleneck Effect}); (3) Mutation; (4) Genetic Recombination during meiosis; (5) Natural Selection.
  \item \textbf{Three Modes of Natural Selection on Phenotypic Traits:}
  \begin{enumerate}
    \item \textbf{Stabilizing Selection:} Favors intermediate phenotype ($2pq$); phenotypic curve becomes narrower and taller (e.g., human infant birth weight centered at $3-3.5\,\mathrm{kg}$).
    \item \textbf{Directional Selection:} Favors one extreme phenotype; shifts population mean towards one direction (e.g., Industrial Melanism in \textit{Biston betularia}, antibiotic resistance in bacteria).
    \item \textbf{Disruptive Selection:} Favors both extreme phenotypes over intermediate; produces a bimodal, two-peaked curve.
  \end{enumerate}
\end{itemize}
\end{formulabox}

\begin{conceptbox}[Chronological Geological Timeline of Human Evolution]
\begin{center}
\scriptsize
\begin{tabular}{|p{2.2cm}|p{1.6cm}|p{2.0cm}|p{1.5cm}|p{6.8cm}|}
\hline
\textbf{Ancestor Stage} & \textbf{Geological Era} & \textbf{Fossil Provenance} & \textbf{Cranial Capacity} & \textbf{Anatomical \& Cultural Characteristics} \\
\hline
\textit{Dryopithecus} & $\approx 15\,\mathrm{mya}$ (Miocene) & Africa and Eurasia & $\approx 350-400\,\mathrm{cc}$ & Arboreal, hairy, arms longer than legs; walked quadrupedally like modern gorillas and chimpanzees. \\
\hline
\textit{Ramapithecus} & $\approx 15\,\mathrm{mya}$ (Miocene) & Shivalik hills (India) & $\approx 400-450\,\mathrm{cc}$ & More man-like than \textit{Dryopithecus}; upright posture tendency; walked semi-bipedally. \\
\hline
\textit{Australopithecus} & $\approx 2-3\,\mathrm{mya}$ (Pliocene) & East African grasslands & $\approx 450-600\,\mathrm{cc}$ & \textbf{Bipedal locomotion}, upright posture; hunted with stone weapons but \textbf{essentially ate fruit}. \\
\hline
\textit{Homo habilis} & $\approx 1.7-2\,\mathrm{mya}$ (Pleistocene) & Olduvai Gorge (Tanzania) & $\mathbf{650-800\,\mathrm{cc}}$ & \textbf{First human-like hominid} (``Tool maker'' / Handy man); constructed crude stone tools; \textbf{did not eat meat}. \\
\hline
\textit{Homo erectus} & $\approx 1.5\,\mathrm{mya}$ (Pleistocene) & Java (1891 fossils) & $\mathbf{900\,\mathrm{cc}}$ & Fully erect posture; proficient toolmaker; first to master \textbf{fire}; \textbf{ate meat}. \\
\hline
\textit{Homo neanderthalensis} & $100,000-40,000$ ya & Neander Valley (Germany) & $\mathbf{1400\,\mathrm{cc}}$ & Heavy brow ridges, receded chin; \textbf{used animal hides to protect bodies; buried their dead with rituals}. \\
\hline
\textit{Homo sapiens} & $75,000-10,000$ ya & Emerged in Africa & $\mathbf{1350-1450\,\mathrm{cc}}$ & Modern man; migrated across continents; prehistoric cave art ($\approx 18,000$ ya at Bhimbetka); \textbf{agriculture began $\approx 10,000$ ya}. \\
\hline
\end{tabular}
\end{center}
\end{conceptbox}

\begin{trapbox}[High-Frequency NEET Assertion-Reason Traps in Health \& Evolution]
\begin{itemize}[leftmargin=*]
  \item \textbf{Rhino Virus vs Streptococcus:} Rhino viruses infect the \textbf{upper respiratory tract and nasal passage, strictly sparing the pulmonary alveoli}; \textit{Streptococcus pneumoniae} infects the alveoli causing fluid accumulation and cyanosis.
  \item \textbf{Malarial Host Nomenclature:} Sexual reproduction occurs \textbf{only in the female \textit{Anopheles} mosquito} (the definitive/primary host). Asexual multiple fission occurs in humans (intermediate host).
  \item \textbf{Morphine vs Heroin:} Morphine is a natural alkaloid; Heroin is a \textbf{semi-synthetic derivative produced by acetylation of morphine}.
  \item \textbf{Darwin vs de Vries on Variation:} Darwinian variations are \textbf{small, continuous, and directional}; de Vriesian mutations are \textbf{single-step large mutations (Saltation), random, and directionless}.
  \item \textbf{Meat Consumption in Hominids:} \textit{Homo habilis} \textbf{did not eat meat}; \textit{Homo erectus} \textbf{ate meat}!
\end{itemize}
\end{trapbox}

\end{document}
